\section{Erd\H{o}s Problem \#521: independent continuation}
\noindent Date: 2026-09-23. Research attempt; no claim of a new published result.
\subsection{1) TARGET AND SOURCE AUDIT}
For the coupled random polynomials $f_n(x)=\sum_{j=0}^n\epsilon_jx^j$ with independent random signs, does the real-root count satisfy $R_n/\log n\to2/\pi$ almost surely along every degree?

All supplied versions considered: \texttt{521.tex}, \texttt{521\_v2.tex}.
Version 1 proves a deterministic subsequence consequence of one-point concentration. Version 2 proves convergence outside a random set of finite reciprocal sum, but labels the result only as logarithmic-density-one. In fact it implies ordinary density-one convergence; I prove that upgrade and its limit.
\subsection{2) WORK}
\paragraph{A deterministic strengthening.}
If $B\subseteq\mathbb N$ and $\sum_{n\in B}1/n<\infty$, then $B$ has natural density zero. For fixed $M<N$,
\[
{|B\cap[1,N]|\over N}\le {M\over N}+\sum_{\substack{n\in B\\n>M}}{1\over n}.
\]
First let $N\to\infty$ and then $M\to\infty$. This proves the claim.
Thus the conclusion of version 2 already implies: almost surely there is a random set $G$ of natural density one along which all its stated normalized root counts converge. Each fixed-tolerance exceptional set has natural density zero, so this is also almost sure statistical convergence in the ordinary counting sense.

For completeness, the supplied concentration hypothesis
$\mathbb P(|R_n/\log n-2/\pi|>\epsilon)\le C_\epsilon e^{-c_\epsilon\sqrt{\log n}}$
implies finite exceptional reciprocal sum by Tonelli, since substitution $u=\sqrt{\log x}$ gives
$\int_3^\infty e^{-c\sqrt{\log x}}dx/x=\int_{\sqrt{\log3}}^\infty2ue^{-cu}du<\infty$.
A diagonal choice of cutoffs for tolerances $1/j$ gives a single $G$: choose each tail reciprocal mass below $2^{-j}$, omit the corresponding bad points in successive cutoff blocks, and also allow the finitely many initial points. The initial finite set was omitted from one displayed estimate in version 2, but it has no effect on finiteness.

This still cannot give convergence at every degree: a deterministic sequence equal to $c+1$ at $n=2^j$ and $c$ otherwise has a harmonic-summable exceptional set and fails to converge.
\subsection{3) ADVERSARIAL VERIFICATION}
No independence between degrees is used by Tonelli. The statistical statement is a consequence of the supplied concentration input, not a new proof of that input. Root counts are not monotone in degree, so density-one convergence cannot be interpolated to all degrees. Reciprocal symmetry gives equal marginal laws, not equality of the coupled stochastic processes.
\subsection{4) FINAL}
\textbf{UNRESOLVED.}
The full almost sure limit still requires control across exceptional degrees, especially outside $[-1,1]$. The stronger density interpretation narrows no such pathwise oscillation by itself.
\subsection{5) SOURCES CHECKED}
Both versions read in full. Checked Do, \texttt{https://arxiv.org/abs/2403.06353}, and Can--Nguyen, \texttt{https://arxiv.org/abs/2311.15446}, on 2026-09-23. Do\textquotesingle s revised abstract explicitly restricts the strong law to $[-1,1]$ and warns that reciprocal symmetry does not transfer it. The detailed concentration inequality is retained as an explicit input, not reproved here.
\subsection{6) COMPLETION ESTIMATE}
This is a conservative subjective measure of progress toward the original question, not a probability of correctness or a claim that the remaining work is routine.
\textbf{COMPLETION: 15\%}
